% =====================================================================
% Section 11: Conclusions and Future Work
% Parent declares \section{Conclusions and Future Work}; this is the body.
% All statistics from _DRAFTING_BRIEF.md (authoritative).
% =====================================================================

\subsection{Conclusions}

This paper has extended an open-source rocket trajectory simulator from a
subsonically valid Barrowman baseline~\cite{barrowman1967,niskanen2009} to a
compressible framework whose integrated apogee accuracy is quantified, and
benchmarked against an independent semi-empirical predictor, across the
supersonic regime to Mach~4.33. The work makes three contributions, framed here
as honestly as the supporting evidence permits.

First, we introduced a \emph{shock-geometry pre-pass} architecture. Once per
timestep the pre-pass walks the airframe from nose to tail and distributes
locally corrected post-shock Mach number, pressure, and temperature to every
downstream drag, normal-force, and stability calculator. The pre-pass is inert
below Mach~1, so it adds no cost to subsonic flight, and its surface-flow
solution is verified bit-for-bit against the classical references: cone
surface conditions agree with Taylor--Maccoll to within a $0.825\%$ shock-angle
error against the gate of $1\%$~\cite{taylormaccoll1933}, and shoulder
expansions reproduce Prandtl--Meyer turning angles to within
$0.004^{\circ}$~\cite{naca1135}. We are deliberate about what this buys. A
mechanism ablation (Sec.~\ref{subsec:prepass-ablation}) shows the pre-pass moves
\emph{integrated apogee} by only $0.15$~percentage points in the mean, because
apogee integrates a trajectory dominated by lower-Mach drag and the pre-pass is
inert below Mach~1. Its value is therefore \emph{local-flow fidelity}---correct
post-shock conditions for fin loads and stability---and its role as the
architectural seam that makes the downstream supersonic models physically
consistent, not a gross-apogee improvement. The dominant apogee mechanism in
that same ablation is the externally calibrated finned-base drag augmentation
($8.10$~percentage points mean effect).

Second, we replaced the underlying engineering models with externally
benchmarked subsystems (Secs.~\ref{sec:drag}--\ref{sec:benchmarks}), each
verified against published wind-tunnel, free-flight range, or computational
data with a documented error metric. These include Van~Driest~II compressible
skin friction~\cite{hopkins1971}, the DATCOM fin wave-drag
method~\cite{datcom1978}, a supersonic base-drag stack validated against
NACA~TN~3393~\cite{chapman1955} and consistent with ESDU~77021~\cite{esdu77021},
and Modified Newtonian impact pressure~\cite{anderson2006}. Representative
component accuracies are total drag on the Basic Finner range
projectile at MAPE~$11.8\%$ (eight points, Mach~1.08--4.30)~\cite{dupuis1997}
and hypersonic cone foredrag at MAPE~$19.7\%$ (eleven points, Mach~6.5--17.2,
maximum~$+57\%$)~\cite{grabow1965}. The verification rests on documented,
re-runnable benchmarks rather than on the author's own computational fluid
dynamics; the comparators are published third-party datasets.

Third, we validated the integrated simulator against a reproducible 25-flight
ground-truth corpus spanning Mach~0.54 to~4.33. The corpus inclusion rule is
external: the flights are taken from Rogers' public RASAero~II
altitude-comparison set~\cite{rogers2015,rogers_rasaero_alt} and comprise
twenty-three single-stage flights and two two-stage flights (AeroPac~104K at
Mach~3.04 and MESOS~293K at Mach~4.33), so the accuracy statistics are not
outcome-curated by the present author. The mean signed apogee error is
$-0.38\%$ (95\% bootstrap confidence interval $[-2.41,\,+1.72]$, 20\,000
resamples), with standard deviation $5.44\%$, RMSE~$5.34\%$, MAE~$4.74\%$, and
all 25~flights within $\pm10\%$ of measured apogee (14 within $\pm5\%$). The
mean-error interval brackets zero, so the predictor is statistically unbiased on
this corpus, and the bias-squared fraction of mean-square error is $0.01$,
indicating that the residual is essentially per-flight random scatter---build
tolerance, motor-lot variation, and atmospheric soundings---rather than
directional model bias. On the 25 paired flights the difference in absolute
error against RASAero~II is $-0.60$~percentage points (95\% bootstrap CI
$[-2.16,\,+0.96]$, straddling zero), and a Wilcoxon signed-rank test finds no
significant difference ($W=143.0$, $p=0.615$)~\cite{wilcoxon1945}. The honest
claim is therefore \emph{parity} with this version-locked RASAero~II set, not
superiority; the recorded RASAero~II values are Rogers' archived predictions
rather than independently rerun pre-flight cases. Because two base-drag scale
constants are corpus-frozen, the headline is partly in-sample; a decontaminated
prospective holdout that places every flight those constants touched into the
development split gives a development MAE of $5.47\%$ (n$=13$) and a genuinely
blind holdout MAE of $3.95\%$ (n$=12$). The holdout being \emph{more} accurate
than development is direct evidence that the two constants generalize rather
than overfit.

\subsection{Future Work}

Concrete next steps, in priority order, follow from the disclosed limitations of
Sec.~\ref{sec:limitations}.

\begin{enumerate}
  \item \textbf{Close the transonic blending bias.} The transonic regime
    (Mach~0.8--1.3) carries the largest per-regime signed error, $-3.66\%$,
    indicating that component wave-drag contributions are presently underblended
    through the sonic line. Redistributing them through a Whitcomb
    equivalent-body-of-revolution area rule, and re-anchoring the blend against
    transonic range data, directly targets this bias.
  \item \textbf{High-Mach validation with non-reconstructed vehicles.} The
    headline corpus reaches only Mach~4.33; the hypersonic regime is empty, and
    the exploratory high-Mach sounding-rocket results
    (Sec.~\ref{sec:exploratory}) are dominated by motor- and
    geometry-reconstruction uncertainty on poorly documented historical flights.
    Admitting sounding rockets on merit requires ground-truth vehicles with
    independently documented thrust curves and geometry, not reconstructions, so
    that an above-Mach-5 integrated claim rests on measurement rather than on
    inferred inputs.
  \item \textbf{Recover the MESOS regression.} The two-stage MESOS~293K closure
    predicts $-6.96\%$ (273{,}056~ft) from the current archived code release---a
    disclosed regression relative to the earlier RFD~v1.2 snapshot ($291{,}601$~ft,
    $-0.6\%$), which was produced by an earlier code version. The $-6.96\%$ value
    is itself reproducible from the present release, and the flight remains within
    the $\pm10\%$ band; the root cause is under investigation, and the next RFD
    release will re-sync the model column. Isolating and recovering this regression
    is required before the two-stage staging-and-coast model can be presented as
    fully validated.
  \item \textbf{Anchor the pitch-damping multiplier to wind-tunnel data.} The
    transonic pitch-damping ($C_{m_q}$) augmentation relies on a single $3\times$
    multiplier that overshoots reference computations by 110--160\% per
    Sznajder~\cite{sznajder2025} and is held at B-level. Lowering it to a
    wind-tunnel-anchored value, validated at flight scale, would raise this
    subsystem to A-level confidence.
  \item \textbf{Ship closed-loop CFD comparators.} The CFD validation
    (Sec.~\ref{sec:cfd}) cites four published third-party datasets rather than
    runs the author can drive end to end. Building an own AGARD-B model file to
    drive the Vidanovi\'c SST $k$--$\omega$ reference~\cite{vidanovic2014} and an
    own RM-10-class ogive--cylinder--boattail file to drive the Sahu base-drag
    reference~\cite{sahu1983} would remove the ``reference dataset only'' caveat.
\end{enumerate}

All source code, the 25-flight Rocket Flight Database~\cite{rfd_zenodo}, and the
analysis scripts that regenerate every figure and table are released under
permissive licenses for end-to-end reproduction. Community contributions on any
of the above items---particularly the transonic blending fix and the MESOS
regression---are warmly invited through the open-source release.
